\documentclass[10pt,a4paper]{amsart}
\usepackage[margin=19mm]{geometry}
\usepackage{amsmath,amssymb,mathtools}
\usepackage[scr=rsfs,cal=euler]{mathalfa}
\usepackage{xfrac}
\usepackage[unicode,hidelinks,bookmarks=false]{hyperref}
\newcommand{\ie}{i.e.\ }
\newcommand{\cf}{cf.\ }
\newcommand{\resp}{respectively\ }
\newcommand{\Subsection}{\subsection}
\providecommand{\nobreakdash}{\nobreak\hbox{-}}
\setlength{\emergencystretch}{2em}
\multlinegap0pt
\usepackage{tikz}
\newtheorem{theorem}{Theorem}
\DeclareMathOperator{\id}{id}
\DeclareMathOperator{\lab}{lab}
\title{Graph quandles: Generalized Cayley graphs of racks and right quasigroups}
\author{Luc Ta}
\date{Source: arXiv:2506.04437v5 (CC BY 4.0)}
\begin{document}
\maketitle

\noindent\emph{Source and licence.} Graph quandles: Generalized Cayley graphs of racks and right quasigroups. Version arXiv:2506.04437v5 (CC BY 4.0), distributed by arXiv under the Creative Commons Attribution 4.0 International licence, http://creativecommons.org/licenses/by/4.0/, as recorded in the arXiv licence metadata for this exact version and shown on the abstract page https://arxiv.org/abs/2506.04437v5. Full licence notice: \texttt{licenses/arxiv-2506.04437-CC-BY-4.0.txt}.\par\smallskip

\noindent\textit{Editorial scope.} Counted unit: the article's theorem thm:lab2 (source0.tex lines 1053--1060). The article's complete original proof of it is delivered with it (source0.tex lines 1062--1090, one \texttt{proof} environment of 29 lines). No mathematics is written, repaired or re-derived: the statement and the proof are the article's own lines, in the article's own notation, and no retained line is substituted or re-flowed.  Retained uncounted as the source-local context the proof needs: lines 191--193; lines 943--945.  Nothing was omitted from any retained span, so the package carries no omission record.  No retained line contains a citation, so the wrapper carries no bibliography and no bibliography entry.  No retained line contains a figure environment, an image-inclusion command or drawing code, so no image is delivered and the package declares no asset dependency.  Editorial formatting only, in addition: an AMS \texttt{amsart} preamble that restates the article's own declarations and macros used by the retained lines, namely the environments \texttt{theorem} on separate counters, together with the standard packages those lines need.  The retained environments print in this document as Theorem 1 in order of appearance, whereas the article prints the same items with the numbering of its own full document.  The complete native source of the article as distributed in the arXiv e-print for this exact version is bundled unchanged as \texttt{sources/179339/source0.tex}.
			\begin{equation}\label{eq:rack}
				R_v R_w = R_{R_v (w)}R_v
			\end{equation}

	\begin{theorem}\label{thm:lab1}
		A labeled digraph is the labeled Cayley digraph of a~right-cancellative magma \textup{(}resp.\ right-divisible magma, right quasigroup\textup{)} if and only if it is an element of $\mathcal{D}$ that is codeterministic \textup{(}resp.\ target-complete, contained in $\mathcal{Q}$\textup{)}.
	\end{theorem}

	\begin{theorem}\label{thm:lab2}
		Let $\Gamma=(V,E,L)$ be a~labeled digraph.
		\begin{enumerate}
			\item $\Gamma$ is the labeled Cayley digraph of a~rack if and only if $\Gamma$ lies in $ \mathcal{Q}$ and satisfies the two rack conditions.
			\item $\Gamma$ is the labeled Cayley digraph of a~quandle if and only if $\Gamma$ lies in $ \mathcal{Q}$, satisfies the two rack conditions, and is label-idempotent.
			\item $\Gamma$ is the labeled Cayley digraph of an involutory right quasigroup if and only if $\Gamma$ is a~label-involutory element of $ \mathcal{Q}$.
		\end{enumerate}
	\end{theorem}

	\begin{proof}
		By Theorem~\ref{thm:lab1}, we can assume that $\Gamma\in\mathcal{Q}$. As noted before, this inclusion implies that $\Gamma=\Gamma_{\lab}(V^\Gamma_R,L)$.

		(1): First, suppose that $V^\Gamma_R$ is a~rack; we show that $\Gamma$ satisfies the rack conditions. For all vertices $v\in V$ and edges ${(v,\ell_i,w_i )\in E}$, we have $w_i =R_{\ell_i}(v)$. It follows from Equation~\eqref{eq:rack} that
		\[
			R_{\ell_1}(w_2)=R_{\ell_1}R_{\ell_2}(v)=R_{R_{\ell_1}(\ell_2)}R_{\ell_1}(v)=R_{R_{\ell_1}(\ell_2)}(w_1).
		\]
		Hence, $\Gamma$ satisfies the first rack condition. Next, let $(v,\ell,w)\in E$, and let $x\in V\setminus L$ satisfy $R_\ell(x)\in L$. Since $R_x =\id_V$ and $R_\ell(v)=w$, Equation~\eqref{eq:rack} yields
		\[
			R_{R_\ell(x)}(w)
            = R_{R_\ell(x)}R_\ell(v)
            = R_\ell R_x(v)
            = R_\ell(v)
            = w.
		\]
		Since $\Gamma=\Gamma_{\lab}(V^\Gamma_R,L)$, it follows that $(w,R_{\ell}(x),w)\in E$, so $\Gamma$ satisfies the second rack condition.

		Conversely, suppose that $\Gamma\in \mathcal{Q}$ satisfies the two rack conditions; we show that $V^\Gamma_R$ is a~rack. We have to verify that
		\begin{equation*}
			R_{\ell_1}R_{\ell_2}(v)=R_{R_{\ell_1}(\ell_2 )}R_{\ell_1}(v)
		\end{equation*}
		for all $\ell_1,\ell_2,v\in V$. If $\ell_1\notin L$, we are done. Next, suppose that $\ell_1\in L$ and $\ell_2\notin L$. If $R_{\ell_1}(\ell_2 )\notin L$, we are done. Otherwise, applying the second rack condition to the edge $(v,\ell_1,R_{\ell_1}(v))$ yields the desired equality. Finally, if $\ell_1,\ell_2\in L$, then $(v,\ell_i,w_i )\in E$ with $w_i:=R_{\ell_i}(v)$. Since $\Gamma$ satisfies the first rack property,
		\[
			R_{\ell_1}R_{\ell_2}(v)=R_{\ell_1}(w_2)=R_{R_{\ell_1}(\ell_2)}(w_1)=R_{R_{\ell_1}(\ell_2)}R_{\ell_1}(v).
		\]
		(2): By the previous claim, it suffices to show that $V^\Gamma_R$ is a~quandle if and only if $\Gamma$ satifies the label-idempotence condition. But this is clear from the construction of $V^\Gamma_R$.

		(3): Clear from the construction of $V^\Gamma_R$.
	\end{proof}

\end{document}
